% !TEX root = ../../main.tex

\section{Discusión}
\label{sec:resultados-discusion}

Respecto de \textbf{PV1}, los 93 programas sometidos al flujo de validación de la solución mostrado en la Figura~\ref{fig:validacion-flujo} conservaron el código de salida y la salida observable en \texttt{ApplicativeDo}, en la selección automática y en cada reordenamiento válido forzado. La coincidencia abarcó las familias estructurales de \texttt{Maybe} y \texttt{Dist.T Rational}, además de pesos no uniformes, soportes dependientes y condicionamiento probabilístico. Los casos analizados en la Sección~\ref{sec:resultados-casos-destacables} muestran que esta preservación se mantuvo ante grafos y construcciones sintácticas diferentes. La evidencia permite responder favorablemente la pregunta para los programas estudiados bajo la hipótesis explícita de conmutatividad, pero continúa siendo observacional y no constituye una demostración de equivalencia semántica para todo programa posible.

Respecto de \textbf{PV2}, en los mismos 93 programas la selección automática coincidió con el primer mínimo recalculado desde los costos de los candidatos compilados individualmente. La Tabla~\ref{tab:resultados-perfiles-costo} y sus cuatro casos muestran que la política se comportó correctamente cuando todos los candidatos empataron, cuando el orden original no fue mínimo, cuando existió un único mínimo y cuando varios reordenamientos compartieron el menor costo. Que un programa no reduzca el costo de \texttt{ApplicativeDo} no representa un fallo: PV2 exige seleccionar un mínimo entre los planes generados, no producir necesariamente una mejora estricta. Como el orden original forma parte del espacio explorado, el costo elegido tampoco puede superar el plan que \texttt{ApplicativeDo} obtiene sobre ese orden.

La cantidad de mejoras estrictas observadas no permite estimar la utilidad de la extensión en programas reales. Los corpus son sintéticos y la mayoría de sus casos fue diseñada para aislar preservación observable, dependencias, binders o empates, no para construir un orden fuente desfavorable. Las cinco variantes que sí mejoraron fueron replicadas con ambas mónadas y preparadas precisamente para admitir un plan menor. El resultado relevante no es su proporción dentro del corpus, sino que la extensión reconoció la alternativa mínima cuando existía y mantuvo un mínimo cuando no había una alternativa mejor.

Respecto de \textbf{PV3}, los seis controles distinguieron correctamente las rutas estándar y conmutativa. Con \texttt{ApplicativeDo} habilitado, los cuatro bloques \texttt{do} ordinarios mantuvieron en las trazas del Renamer \texttt{commutative-do = False}, aun en presencia de \texttt{QualifiedDo}, la importación de la interfaz y la instancia correspondiente. Los dos bloques \texttt{CD.do} produjeron \texttt{commutative-do = True} y generaron seis candidatos cada uno, uno original y cinco alternativos. Estos resultados confirman la activación local mediante el marcador dentro de las configuraciones evaluadas. La instancia \texttt{CommutativeMonad} sigue siendo necesaria para utilizar las operaciones calificadas, pero expresa una restricción comprobada durante el tipado y no constituye por sí sola el disparador del Renamer.

Los grafos presentados también muestran que el tamaño del espacio de búsqueda y la calidad del plan son propiedades relacionadas, pero diferentes. Una cadena completa produjo un único candidato; un grafo vacío de cuatro vértices produjo 24; y dos cadenas independientes de longitud cuatro produjeron 70. Sin embargo, el diamante mantuvo el mínimo que \texttt{ApplicativeDo} ya alcanzaba, mientras el orden desfavorable de las dos cadenas permitió reducir el costo de 6 a 4. Del mismo modo, las familias de \texttt{let}, binders, distribuciones dependientes, condicionamiento y \textit{shadowing} confirmaron que las aristas se forman a partir de las identidades y lecturas locales detectadas por el Renamer, no solo a partir de binders monádicos simples.

Finalmente, el control con \texttt{IO} de la Sección~\ref{sec:resultados-io} separa la validez estructural de la equivalencia semántica. El candidato invertido preservó el grafo RAW vacío y retornó el mismo entero, pero cambió el orden de los efectos visibles. Su divergencia no contradice los resultados positivos, sino que confirma el límite previsto por el diseño: las precedencias RAW protegen el flujo local de datos, mientras la posibilidad de intercambiar efectos depende de la conmutatividad declarada. Asimismo, las reducciones informadas corresponden al modelo estructural uniforme y no a mediciones de tiempo de ejecución. Bajo estas condiciones, los resultados responden favorablemente las tres preguntas de validación dentro del alcance delimitado en la Sección~\ref{sec:validacion-alcance}.
